When monochromatic light passes through a diffraction grating, it
splits into a characteristic pattern of bright and dark fringes due to
wave interference.
Diffraction, first described by Francesco Grimaldi in 1660, is the
phenomenon where waves bend and spread around obstacles or through
apertures.
The extent of this bending depends on the ratio of the wavelength to
the size of the opening.

A diffraction grating consists of many parallel slits with spacing
comparable to the wavelength of light.
When light passes through these slits, the emerging waves interfere
with each other, creating a pattern of constructive and destructive
interference.
Constructive interference produces bright fringes (maxima), while
destructive interference creates dark regions (minima).
The central, undeviated beam is the zeroth-order maximum, and
subsequent symmetrically spaced fringes are the higher-order maxima.

\begin{figure}[H]
\centering
\includegraphics[width=0.75\textwidth]{../out/intensity_plots/air_v2_intensity}
\caption{Typical diffraction pattern intensity profile measured through
air using the 650~nm laser.
The central peak (zeroth order) shows maximum intensity with symmetric
first-order fringes on either side, demonstrating the characteristic
diffraction pattern produced by the grating.}
\label{fig:diffraction_pattern_intro}
\end{figure}

Constructive interference occurs when the path difference between waves
from adjacent slits is an integer multiple ($m$) of the wavelength
($\lambda$).
For a diffraction grating with slit spacing $d$, this path difference
forms the opposite side of a right triangle where $d$ is the
hypotenuse.
The angle $\theta$ of the $m^{\text{th}}$-order fringe relative to the
normal is given by:
\[\theta_{air} = \arcsin \left(\frac{m \lambda}{d}\right)\tag{1}\]

The refractive index of a medium is defined as the ratio of the speed
of light in a vacuum ($c$) to the speed of light in that medium ($v$).
For linear media, this index depends on optical density and can be
affected by factors such as temperature and pressure.
Since the frequency ($f$) of light remains constant as it propagates
through different media (being determined by the source), we can relate
the wavelength in vacuum ($\lambda_0$) to the wavelength in a medium
($\lambda$).

Starting with the fundamental relationship $c = f \lambda_0$ for light
in vacuum and $v = f \lambda$ for light in a medium, we can solve for
frequency and equate them:
\[\frac{c}{\lambda_0} = \frac{v}{\lambda}\tag{2}\]

This yields the relationship between wavelength and refractive index:
\[\lambda = \frac{\lambda_0}{n} \quad \text{where} \quad n = \frac{c}{v}\tag{3}\]

Combining this with the grating equation allows us to determine the
refractive index from the diffraction angle:
\[n = \frac{\lambda_0 m}{d \sin(\theta)}\tag{4}\]

This formula is valid when light travels through a single uniform
medium.
However, in this experiment, the diffraction grating remains in air
while the diffracted light passes through a liquid-filled cuvette
before reaching the detector.
This multi-medium optical path requires a more detailed analysis.

Since the grating cannot be submerged, it remains in air
($n_{air} \approx 1$), and equation (1) determines the initial
diffraction angle in air.
The light then enters the liquid medium, where it refracts according to
Snell's law:
\[n_1 \sin(\theta_1) = n_2 \sin(\theta_2)\tag{5}\]

As the beam passes through the liquid, refraction bends it toward the
normal if the liquid is more optically dense than air, or away from the
normal if less dense.
This refraction causes a lateral displacement ($\Delta x$) of the
fringe position at the detector.
For accurate measurements, any air gap greater than 1~mm between the
grating and the liquid must be accounted for, as such gaps introduce
significant systematic errors.

The lateral displacement of a fringe depends on the path length through
each medium and the refractive indices involved.
For a fringe at position $x_i$, the geometry of the beam path gives:
\[x_i = p \tan(\theta_i)\tag{6}\]
where $p$ is the path length and $\theta_i$ is the propagation angle in
that medium.

For the two-medium system (air and liquid) with the beam initially
normal to the grating, the displacement of the $i^{\text{th}}$-order
fringe can be expressed as:
\[\Delta x = p \sin(\theta_{\text{air}}) \left(\frac{1}{\sqrt{1 - \sin^2(\theta_{\text{air}})}} - \frac{1}{\sqrt{n_{liq}^2 - \sin^2(\theta_{\text{air}})}}\right)\tag{7}\]
where $p$ is the liquid path length.
This equation assumes normal incidence on the grating and parallel
cuvette walls, and is not suitable for more than two media or
non-normal incidence without modification.

Solving equation (7) for the refractive index yields:
\[n_{liq} = \sin(\theta_{\text{air}}) \sqrt{1 + \frac{p^2 (1 - \sin^2(\theta_{\text{air}}))}{(p \sin(\theta_{\text{air}}) - \Delta x \sqrt{1 - \sin^2(\theta_{\text{air}})})^2}}\tag{8}\]

By measuring the fringe positions with an empty cuvette (baseline) and
comparing them to the positions with the filled cuvette, the effective
displacement $\Delta x$ can be determined.
This approach eliminates the need to account for the optical effects of
the cuvette walls, as they are present in both measurements and cancel
out when calculating the displacement.

Once the refractive index is determined, the wavelength of light in the
liquid can be calculated by combining equations (3) and (8):
\[\lambda_{liq} = \frac{d}{m \sqrt{1 + \frac{p^2 (1 - \sin^2(\theta_{\text{air}}))}{(p \sin(\theta_{\text{air}}) - \Delta x \sqrt{1 - \sin^2(\theta_{\text{air}})})^2}}}\tag{9}\]

This wavelength shift is most noticeable in materials with high
refractive indices.
The goal of this experiment is to validate this cost-effective
diffraction-based method for measuring the refractive indices of common
liquids.

